% =============================================================================
%  PAES Studio: presentación de repaso (plantilla Beamer 16:9)
%  Lista para Overleaf: Nuevo proyecto → Subir proyecto → compilar con pdfLaTeX.
%
%  Estética minimalista con la paleta Zinc de la app: sin degradados, sin
%  símbolos de navegación, líneas finas y alternativas con su tecla A–E.
%  Comandos:
%    \alternativa{A}{texto}       alternativa
%    \alternativaok{C}{texto}     alternativa correcta (para la diapositiva de solución)
%    \paso{n}{texto}      paso numerado del desarrollo
%    \ejebarra{Eje}{0.42}{42 %}   barra proporcional de frecuencia
%  Entornos: cajaclave, cajatrampa, lineatiempo (filas «fecha & hecho \\»).
%
%  El generador scripts/pipeline/generate_latex_summary.py reemplaza los bloques
%  delimitados por %%<NOMBRE>%% … %%</NOMBRE>%%.
% =============================================================================
\documentclass[aspectratio=169,11pt]{beamer}

\usepackage{lmodern}        % antes de fontenc/babel: no depende de las fuentes EC
\usepackage[T1]{fontenc}
\usepackage[utf8]{inputenc}
\usepackage[spanish,es-noshorthands,es-tabla]{babel}
\spanishplainpercent        % \% sin espaciado automático (evita choques en modo matemático)
\usepackage{textcomp}
\usepackage{amsmath,amssymb}
\usepackage{booktabs}
\usepackage{array}
\usepackage[most]{tcolorbox}

% ── Paleta Zinc ──────────────────────────────────────────────────────────────
\definecolor{zinc950}{HTML}{09090B}
\definecolor{zinc900}{HTML}{18181B}
\definecolor{zinc700}{HTML}{3F3F46}
\definecolor{zinc500}{HTML}{71717A}
\definecolor{zinc300}{HTML}{D4D4D8}
\definecolor{zinc200}{HTML}{E4E4E7}
\definecolor{zinc50}{HTML}{FAFAFA}
\definecolor{esmeralda}{HTML}{047857}
\definecolor{rosa}{HTML}{BE123C}

\newcommand{\sen}{\operatorname{sen}}

% ── Tema mínimo ──────────────────────────────────────────────────────────────
\usefonttheme{professionalfonts}
\setbeamertemplate{navigation symbols}{}
\setbeamercolor{background canvas}{bg=white}
\setbeamercolor{normal text}{fg=zinc900}
\setbeamercolor{structure}{fg=zinc900}
\setbeamercolor{alerted text}{fg=zinc950}
\setbeamercolor{item}{fg=zinc500}
\setbeamerfont{frametitle}{series=\bfseries,size=\large}
\setbeamertemplate{itemize items}{\raisebox{0.15ex}{\color{zinc500}\scriptsize$\bullet$}}
\setbeamertemplate{enumerate items}{\color{zinc500}\insertenumlabel.}
\setbeamersize{text margin left=1.6cm,text margin right=1.6cm}

\makeatletter
\setbeamertemplate{frametitle}{%
  \vspace*{0.9em}%
  {\usebeamerfont{frametitle}\color{zinc950}\insertframetitle}\par
  \ifx\insertframesubtitle\@empty\else{\small\color{zinc500}\insertframesubtitle}\par\fi
  \vspace{0.25em}{\color{zinc200}\rule{\textwidth}{0.4pt}}%
}
\makeatother
\setbeamertemplate{footline}{%
  \hspace*{1.6cm}{\color{zinc500}\tiny\insertshorttitle}\hfill
  {\color{zinc500}\tiny\ttfamily\insertframenumber/\inserttotalframenumber}\hspace*{1.6cm}\par\vspace*{0.7em}%
}
\setbeamertemplate{title page}{%
  \vfill
  {\color{zinc500}\small PAES Studio}\par\vspace{0.6em}
  {\usebeamerfont{title}\bfseries\Huge\color{zinc950}\inserttitle}\par\vspace{0.5em}
  {\large\color{zinc700}\insertsubtitle}\par\vspace{1.2em}
  {\color{zinc300}\rule{0.35\textwidth}{0.4pt}}\par\vspace{0.6em}
  {\small\color{zinc500}\insertdate}\par
  \vfill
}

% ── Componentes ──────────────────────────────────────────────────────────────
\newtcbox{\tecla}{on line, boxrule=0.4pt, arc=1.5pt, colback=zinc50, colframe=zinc300,
  left=3pt, right=3pt, top=1pt, bottom=1pt, fontupper=\ttfamily\footnotesize\color{zinc700}}
\newtcbox{\teclaok}{on line, boxrule=0.4pt, arc=1.5pt, colback=esmeralda, colframe=esmeralda,
  left=3pt, right=3pt, top=1pt, bottom=1pt, fontupper=\ttfamily\footnotesize\bfseries\color{white}}
\newtcolorbox{opcion}[1][]{enhanced, boxrule=0.4pt, arc=2pt, colback=white, colframe=zinc300,
  left=6pt, right=6pt, top=3pt, bottom=3pt, before skip=3pt, after skip=3pt, #1}
\newcommand{\alternativa}[2]{\begin{opcion}\tecla{#1}\hspace{0.6em}#2\end{opcion}}
\newcommand{\alternativaok}[2]{\begin{opcion}[colframe=esmeralda!60, colback=esmeralda!5]\teclaok{#1}\hspace{0.6em}#2\end{opcion}}
\newcommand{\paso}[2]{\par\noindent\makebox[1.4em][l]{\color{zinc500}\ttfamily\small #1}\parbox[t]{\dimexpr\linewidth-1.4em}{#2}\par\vspace{0.45em}}
\newtcolorbox{cajaclave}[1][]{enhanced, boxrule=0.4pt, arc=2pt, colback=zinc50, colframe=zinc300,
  borderline west={2pt}{0pt}{zinc900}, left=7pt, right=7pt, top=3pt, bottom=4pt,
  fonttitle=\bfseries\small, coltitle=zinc900, colbacktitle=zinc50, titlerule=0pt, title={#1}}
\newtcolorbox{cajatrampa}[1][Trampas frecuentes]{enhanced, boxrule=0.4pt, arc=2pt, colback=white, colframe=rosa!35,
  borderline west={2pt}{0pt}{rosa}, left=7pt, right=7pt, top=3pt, bottom=4pt,
  fonttitle=\bfseries\small, coltitle=zinc900, colbacktitle=white, titlerule=0pt, title={#1}}
\newenvironment{lineatiempo}
  {\footnotesize\renewcommand{\arraystretch}{1.15}\begin{tabular}{@{}>{\leavevmode\bfseries\color{zinc700}}p{0.14\textwidth}p{0.74\textwidth}@{}}}
  {\end{tabular}}
\newlength{\anchobarra}
\newcommand{\ejebarra}[3]{%
  \par\noindent\makebox[0.34\linewidth][l]{\small #1}%
  \setlength{\anchobarra}{0.5\linewidth}%
  \makebox[\anchobarra][l]{{\color{zinc200}\rule[0.2ex]{\anchobarra}{1.1ex}}\hspace*{-\anchobarra}{\color{zinc900}\rule[0.2ex]{#2\anchobarra}{1.1ex}}}%
  \hspace{0.8em}{\small\ttfamily #3}\par\vspace{0.5em}}

%%<META>%%
\title[Repaso PAES: Historia y Ciencias Sociales]{Repaso PAES}
\subtitle{Historia y Ciencias Sociales}
\date{1 de noviembre de 2026}
%%</META>%%

\begin{document}

\begin{frame}[plain]
  \titlepage
\end{frame}

%%<CONTENIDO>%%
\begin{frame}{Contenidos}

  \framesubtitle{5 ejes, en el orden del temario}

  \begin{enumerate}\small

    \item \textbf{Historia: mundo, América y Chile}: Revoluciones y siglo XIX en Europa, Chile en el siglo XIX: independencia, república y expansión, Cuestión social y crisis del parlamentarismo, Guerras mundiales y crisis de 1929, Guerra Fría y América Latina, Chile en el siglo XX: del Estado de compromiso a la transición

    \item \textbf{Formación ciudadana}: Estado, Constitución y poderes del Estado, Democracia, derechos humanos y participación

    \item \textbf{Sistema económico}: Problema económico y mercado, Macroeconomía y finanzas personales

    \item \textbf{Geografía y espacio}: Población, territorio y ambiente

    \item \textbf{Habilidades de pensamiento histórico}: Análisis de fuentes y pensamiento histórico

  \end{enumerate}

\end{frame}

\begin{frame}{Qué evalúa la prueba}

  \framesubtitle{Distribución por eje en 7 ensayos (386 preguntas)}

  \vspace{0.8em}

  \ejebarra{Formación ciudadana}{0.438}{44 \%}

  \ejebarra{Historia: mundo, América y Chile}{0.246}{25 \%}

  \ejebarra{Sistema económico}{0.228}{23 \%}

  \ejebarra{Geografía y espacio}{0.088}{9 \%}

\end{frame}

\section{Historia: mundo, América y Chile}

\begin{frame}[plain]
  \vfill
  {\color{zinc500}\small Eje}\par\vspace{0.3em}{\Large\bfseries Historia: mundo, América y Chile}\par\vspace{0.6em}
  {\color{zinc700} Procesos políticos, económicos y sociales de los siglos XIX y XX.}\par\vspace{1em}
  {\small\color{zinc500} Revoluciones y siglo XIX en Europa, Chile en el siglo XIX: independencia, república y expansión, Cuestión social y crisis del parlamentarismo, Guerras mundiales y crisis de 1929, Guerra Fría y América Latina, Chile en el siglo XX: del Estado de compromiso a la transición}
  \vfill
\end{frame}

\begin{frame}[allowframebreaks]{Revoluciones y siglo XIX en Europa}
  \framesubtitle{Teoría esencial}
  \begin{itemize}\small
    \item 1789: fin del Antiguo Régimen en Francia.
    \item Industrialización $\rightarrow$ urbanización y cuestión obrera.
    \item Imperialismo: reparto de África en la Conferencia de Berlín (1884–1885).
  \end{itemize}
  \begin{cajaclave}[Ideas clave del liberalismo]\footnotesize
Soberanía nacional, separación de poderes, derechos individuales (libertad, propiedad, igualdad ante la ley), constitución escrita y libre mercado.
  \end{cajaclave}
\end{frame}

\begin{frame}{Revoluciones y siglo XIX en Europa}
  \framesubtitle{Línea de tiempo}
  \begin{lineatiempo}
  1789 & Revolución Francesa y Declaración de los Derechos del Hombre \\
  c. 1760–1840 & Primera Revolución Industrial en Inglaterra \\
  1815 & Congreso de Viena \\
  1848 & Revoluciones liberales en Europa \\
  1871 & Unificación alemana \\
  1884–1885 & Conferencia de Berlín \\
  \end{lineatiempo}
\end{frame}

\begin{frame}{Desafío}
  \framesubtitle{Revoluciones y siglo XIX en Europa}
  {\small ¿Cuál fue una consecuencia social directa de la Revolución Industrial en Inglaterra?}
  \vspace{0.4em}
  \alternativa{A}{El fortalecimiento de los gremios artesanales.}
  \alternativa{B}{La formación de una clase obrera concentrada en ciudades.}
  \alternativa{C}{La desaparición de la burguesía.}
  \alternativa{D}{El retorno de la población al campo.}
\end{frame}

\begin{frame}[allowframebreaks]{Solución}
  \framesubtitle{Las fábricas concentraron trabajadores asalariados en las ciudades.}
  {\small
  \paso{1}{La producción fabril requería mano de obra cercana.}
  \paso{2}{La migración campo–ciudad formó el proletariado urbano.}
  }
  \alternativaok{B}{La formación de una clase obrera concentrada en ciudades.}
  \begin{cajatrampa}\footnotesize
    \begin{itemize}
    \item Atribuir a la Revolución Francesa el sufragio universal femenino.
    \item Confundir liberalismo político con socialismo.
    \end{itemize}
  \end{cajatrampa}
\end{frame}

\begin{frame}[allowframebreaks]{Chile en el siglo XIX: independencia, república y expansión}
  \framesubtitle{Teoría esencial}
  \begin{itemize}\small
    \item 1810: Primera Junta, en nombre de Fernando VII.
    \item 1833: Constitución de orden conservador y presidencialista.
    \item Guerra del Pacífico $\rightarrow$ riqueza salitrera para el Estado.
  \end{itemize}
  \begin{cajaclave}[Ciclo del salitre]\footnotesize
Exportación de salitre $\rightarrow$ impuestos de exportación $\rightarrow$ grandes ingresos fiscales $\rightarrow$ obras públicas, educación y ferrocarriles. Gran parte de las oficinas eran de capitales extranjeros.
  \end{cajaclave}
\end{frame}

\begin{frame}{Chile en el siglo XIX: independencia, república y expansión}
  \framesubtitle{Línea de tiempo}
  \begin{lineatiempo}
  1810 & Primera Junta Nacional de Gobierno \\
  1814 & Desastre de Rancagua: Reconquista \\
  1818 & Proclamación de la independencia y batalla de Maipú \\
  1833 & Constitución de 1833 \\
  1879–1884 & Guerra del Pacífico \\
  1888 & Incorporación de Rapa Nui \\
  1891 & Guerra Civil \\
  \end{lineatiempo}
\end{frame}

\begin{frame}{Desafío}
  \framesubtitle{Chile en el siglo XIX: independencia, república y expansión}
  {\small ¿Qué hecho desencadenó la formación de juntas de gobierno en América en 1810?}
  \vspace{0.4em}
  \alternativa{A}{La Revolución Industrial.}
  \alternativa{B}{La invasión napoleónica a España y la captura de Fernando VII.}
  \alternativa{C}{La Guerra del Pacífico.}
  \alternativa{D}{La Constitución de 1833.}
\end{frame}

\begin{frame}[allowframebreaks]{Solución}
  \framesubtitle{Sin rey legítimo, surgió la pregunta por quién gobernaba en su nombre.}
  {\small
  \paso{1}{1808: Napoleón invade España y Fernando VII queda cautivo.}
  \paso{2}{Los cabildos americanos forman juntas que gobiernan en nombre del rey.}
  }
  \alternativaok{B}{La invasión napoleónica a España y la captura de Fernando VII.}
  \begin{cajatrampa}\footnotesize
    \begin{itemize}
    \item Creer que la Junta de 1810 declaró la independencia.
    \item Pensar que el Estado explotó directamente las salitreras.
    \end{itemize}
  \end{cajatrampa}
\end{frame}

\begin{frame}[allowframebreaks]{Cuestión social y crisis del parlamentarismo}
  \framesubtitle{Teoría esencial}
  \begin{itemize}\small
    \item 1907: matanza de la Escuela Santa María de Iquique.
    \item 1924: leyes sociales; 1925: nueva Constitución.
    \item Conventillo = símbolo de la vivienda popular urbana.
  \end{itemize}
  \begin{cajaclave}[Respuestas a la cuestión social]\footnotesize
Movimiento obrero (huelgas, organización), Iglesia (encíclica \emph{Rerum Novarum}, 1891), Estado (leyes sociales de 1924) y partidos políticos de base popular.
  \end{cajaclave}
\end{frame}

\begin{frame}{Cuestión social y crisis del parlamentarismo}
  \framesubtitle{Línea de tiempo}
  \begin{lineatiempo}
  1891–1925 & República parlamentaria \\
  1907 & Matanza de la Escuela Santa María de Iquique \\
  1920 & Ley de Instrucción Primaria Obligatoria \\
  1924 & Leyes sociales \\
  1925 & Constitución de 1925 \\
  \end{lineatiempo}
\end{frame}

\begin{frame}{Desafío}
  \framesubtitle{Cuestión social y crisis del parlamentarismo}
  {\small ¿Qué cambio institucional introdujo la Constitución de 1925 en Chile?}
  \vspace{0.4em}
  \alternativa{A}{Estableció el parlamentarismo.}
  \alternativa{B}{Fortaleció el presidencialismo y separó la Iglesia del Estado.}
  \alternativa{C}{Creó la monarquía constitucional.}
  \alternativa{D}{Suprimió el Congreso Nacional.}
\end{frame}

\begin{frame}[allowframebreaks]{Solución}
  \framesubtitle{La Constitución de 1925 respondió a la crisis del parlamentarismo.}
  {\small
  \paso{1}{Devolvió atribuciones al Presidente frente al Congreso.}
  \paso{2}{Estableció la separación entre la Iglesia y el Estado.}
  }
  \alternativaok{B}{Fortaleció el presidencialismo y separó la Iglesia del Estado.}
  \begin{cajatrampa}\footnotesize
    \begin{itemize}
    \item Ubicar la cuestión social en la segunda mitad del siglo XX.
    \item Confundir el período parlamentario con el presidencial.
    \end{itemize}
  \end{cajatrampa}
\end{frame}

\begin{frame}[allowframebreaks]{Guerras mundiales y crisis de 1929}
  \framesubtitle{Teoría esencial}
  \begin{itemize}\small
    \item Causa estructural $\neq$ detonante.
    \item La crisis de 1929 golpeó fuertemente a Chile por la caída de sus exportaciones.
    \item 1945: fin de la guerra y fundación de la ONU.
  \end{itemize}
  \begin{cajaclave}[Multicausalidad de la Primera Guerra Mundial]\footnotesize
Causas estructurales: imperialismo, nacionalismo, paz armada, alianzas. Detonante: asesinato del archiduque Francisco Fernando en Sarajevo (1914).
  \end{cajaclave}
\end{frame}

\begin{frame}{Guerras mundiales y crisis de 1929}
  \framesubtitle{Línea de tiempo}
  \begin{lineatiempo}
  1914–1918 & Primera Guerra Mundial \\
  1917 & Revolución Rusa \\
  1919 & Tratado de Versalles \\
  1929 & Caída de la bolsa de Nueva York \\
  1939–1945 & Segunda Guerra Mundial \\
  1945 & Creación de la ONU \\
  \end{lineatiempo}
\end{frame}

\begin{frame}{Desafío}
  \framesubtitle{Guerras mundiales y crisis de 1929}
  {\small ¿Qué efecto tuvo la crisis de 1929 en la economía chilena?}
  \vspace{0.4em}
  \alternativa{A}{Un auge de las exportaciones de salitre.}
  \alternativa{B}{Una fuerte caída de las exportaciones y alto desempleo.}
  \alternativa{C}{La inmediata industrialización del país.}
  \alternativa{D}{No tuvo efectos por el aislamiento de Chile.}
\end{frame}

\begin{frame}[allowframebreaks]{Solución}
  \framesubtitle{Chile dependía de exportar materias primas, cuya demanda se desplomó.}
  {\small
  \paso{1}{La contracción mundial redujo la demanda de cobre y salitre.}
  \paso{2}{Cayeron los ingresos y aumentó el desempleo.}
  }
  \alternativaok{B}{Una fuerte caída de las exportaciones y alto desempleo.}
  \begin{cajatrampa}\footnotesize
    \begin{itemize}
    \item Presentar el detonante como la única causa de la guerra.
    \item Confundir los bandos de ambas guerras.
    \end{itemize}
  \end{cajatrampa}
\end{frame}

\begin{frame}[allowframebreaks]{Guerra Fría y América Latina}
  \framesubtitle{Teoría esencial}
  \begin{itemize}\small
    \item Guerra «fría»: sin enfrentamiento directo entre superpotencias.
    \item 1962: crisis de los misiles en Cuba.
    \item 1989: caída del Muro de Berlín.
  \end{itemize}
  \begin{cajaclave}[Bloques]\footnotesize
Occidental: Estados Unidos, OTAN (1949), Plan Marshall. Oriental: URSS, Pacto de Varsovia (1955). Fin: caída del Muro de Berlín (1989) y disolución de la URSS (1991).
  \end{cajaclave}
\end{frame}

\begin{frame}{Guerra Fría y América Latina}
  \framesubtitle{Línea de tiempo}
  \begin{lineatiempo}
  1947 & Doctrina Truman y Plan Marshall \\
  1949 & Creación de la OTAN \\
  1959 & Revolución Cubana \\
  1961 & Construcción del Muro de Berlín \\
  1962 & Crisis de los misiles \\
  1989 & Caída del Muro de Berlín \\
  1991 & Disolución de la URSS \\
  \end{lineatiempo}
\end{frame}

\begin{frame}{Desafío}
  \framesubtitle{Guerra Fría y América Latina}
  {\small ¿Qué característica define a la Guerra Fría?}
  \vspace{0.4em}
  \alternativa{A}{Un enfrentamiento armado directo entre EE. UU. y la URSS.}
  \alternativa{B}{La rivalidad entre dos bloques que evitaron una guerra directa.}
  \alternativa{C}{La alianza entre EE. UU. y la URSS contra el fascismo.}
  \alternativa{D}{El conflicto entre Francia e Inglaterra por colonias.}
\end{frame}

\begin{frame}[allowframebreaks]{Solución}
  \framesubtitle{La disuasión nuclear evitó la guerra directa.}
  {\small
  \paso{1}{El conflicto se libró con propaganda, carreras armamentista y espacial y guerras indirectas.}
  \paso{2}{La alternativa C describe la Segunda Guerra Mundial.}
  }
  \alternativaok{B}{La rivalidad entre dos bloques que evitaron una guerra directa.}
  \begin{cajatrampa}\footnotesize
    \begin{itemize}
    \item Describir la Guerra Fría como un enfrentamiento militar directo entre EE. UU. y la URSS.
    \end{itemize}
  \end{cajatrampa}
\end{frame}

\begin{frame}[allowframebreaks]{Chile en el siglo XX: del Estado de compromiso a la transición}
  \framesubtitle{Teoría esencial}
  \begin{itemize}\small
    \item 1939: CORFO, símbolo del Estado empresario.
    \item 1971: nacionalización del cobre, aprobada por unanimidad en el Congreso.
    \item 1988: triunfo del «No» en el plebiscito.
  \end{itemize}
  \begin{cajaclave}[Modelos económicos]\footnotesize
ISI (1938–1973): Estado empresario y protección arancelaria. Neoliberal (desde 1975): apertura comercial, privatizaciones y rol subsidiario del Estado.
  \end{cajaclave}
\end{frame}

\begin{frame}{Chile en el siglo XX: del Estado de compromiso a la transición}
  \framesubtitle{Línea de tiempo}
  \begin{lineatiempo}
  1938 & Frente Popular: gobierno de Pedro Aguirre Cerda \\
  1939 & Creación de la CORFO \\
  1967 & Ley de reforma agraria \\
  1971 & Nacionalización del cobre \\
  1973 & Golpe de Estado \\
  1980 & Constitución de 1980 \\
  1988 & Plebiscito: triunfo del «No» \\
  1990 & Inicio de la transición: gobierno de Patricio Aylwin \\
  \end{lineatiempo}
\end{frame}

\begin{frame}{Desafío}
  \framesubtitle{Chile en el siglo XX: del Estado de compromiso a la transición}
  {\small ¿Qué política económica caracterizó a Chile entre 1938 y comienzos de los años setenta?}
  \vspace{0.4em}
  \alternativa{A}{La apertura comercial irrestricta.}
  \alternativa{B}{La industrialización por sustitución de importaciones con un Estado empresario.}
  \alternativa{C}{La privatización de las empresas estatales.}
  \alternativa{D}{La economía basada exclusivamente en el salitre.}
\end{frame}

\begin{frame}[allowframebreaks]{Solución}
  \framesubtitle{Tras la crisis de 1929, el Estado impulsó la industria nacional.}
  {\small
  \paso{1}{La CORFO (1939) creó y financió industrias.}
  \paso{2}{Se protegió la producción nacional con aranceles.}
  }
  \alternativaok{B}{La industrialización por sustitución de importaciones con un Estado empresario.}
  \begin{cajatrampa}\footnotesize
    \begin{itemize}
    \item Confundir chilenización (1966) con nacionalización (1971) del cobre.
    \item Atribuir la Constitución de 1980 a la transición.
    \end{itemize}
  \end{cajatrampa}
\end{frame}

\section{Formación ciudadana}

\begin{frame}[plain]
  \vfill
  {\color{zinc500}\small Eje}\par\vspace{0.3em}{\Large\bfseries Formación ciudadana}\par\vspace{0.6em}
  {\color{zinc700} Estado, Constitución, democracia, derechos humanos y participación.}\par\vspace{1em}
  {\small\color{zinc500} Estado, Constitución y poderes del Estado, Democracia, derechos humanos y participación}
  \vfill
\end{frame}

\begin{frame}[allowframebreaks]{Estado, Constitución y poderes del Estado}
  \framesubtitle{Teoría esencial}
  \begin{itemize}\small
    \item Fiscalizar los actos del Gobierno es atribución exclusiva de la Cámara.
    \item La Contraloría controla la legalidad de los actos de la administración.
    \item El Tribunal Constitucional vela por que las leyes respeten la Constitución.
  \end{itemize}
  \begin{cajaclave}[Formación de la ley]\footnotesize
Iniciativa (mensaje del Presidente o moción parlamentaria) $\rightarrow$ discusión y aprobación en ambas cámaras $\rightarrow$ sanción (o veto) del Presidente $\rightarrow$ promulgación $\rightarrow$ publicación en el Diario Oficial.
  \end{cajaclave}
\begin{center}\footnotesize
\begin{tabular}{@{}>{\raggedright\arraybackslash}p{\dimexpr(\linewidth-4\tabcolsep)/3\relax}>{\raggedright\arraybackslash}p{\dimexpr(\linewidth-4\tabcolsep)/3\relax}>{\raggedright\arraybackslash}p{\dimexpr(\linewidth-4\tabcolsep)/3\relax}@{}}
\toprule
\textbf{Poder} & \textbf{Órgano} & \textbf{Función principal} \\
\midrule
Ejecutivo & Presidente de la República & gobernar y administrar el Estado \\
Legislativo & Congreso Nacional & legislar; la Cámara fiscaliza \\
Judicial & Tribunales de justicia & resolver conflictos aplicando la ley \\
\bottomrule
\end{tabular}
\end{center}
\end{frame}

\begin{frame}{Desafío}
  \framesubtitle{Estado, Constitución y poderes del Estado}
  {\small ¿Qué órgano controla la legalidad de los actos de la administración del Estado en Chile?}
  \vspace{0.4em}
  \alternativa{A}{El Senado.}
  \alternativa{B}{La Contraloría General de la República.}
  \alternativa{C}{El Banco Central.}
  \alternativa{D}{El Servicio Electoral.}
\end{frame}

\begin{frame}[allowframebreaks]{Solución}
  \framesubtitle{La Contraloría es un órgano autónomo de control.}
  {\small
  \paso{1}{Revisa que los actos de la administración se ajusten a la ley (toma de razón).}
  \paso{2}{No legisla ni juzga: controla legalidad.}
  }
  \alternativaok{B}{La Contraloría General de la República.}
  \begin{cajatrampa}\footnotesize
    \begin{itemize}
    \item Atribuir al Senado la fiscalización del Gobierno.
    \item Confundir el Tribunal Constitucional con la Corte Suprema.
    \end{itemize}
  \end{cajatrampa}
\end{frame}

\begin{frame}[allowframebreaks]{Democracia, derechos humanos y participación}
  \framesubtitle{Teoría esencial}
  \begin{itemize}\small
    \item Estado de derecho: todos, incluidas las autoridades, se someten a la ley.
    \item Universalidad: los DD. HH. corresponden a toda persona.
    \item Plebiscito = mecanismo de democracia directa.
  \end{itemize}
  \begin{cajaclave}[Democracia directa y representativa]\footnotesize
Representativa: la ciudadanía elige autoridades que deciden. Directa: la ciudadanía decide por sí misma (plebiscito, iniciativa popular, consulta).
  \end{cajaclave}
\end{frame}

\begin{frame}{Desafío}
  \framesubtitle{Democracia, derechos humanos y participación}
  {\small Que los derechos humanos sean inalienables significa que:}
  \vspace{0.4em}
  \alternativa{A}{Solo los reconocen algunos países.}
  \alternativa{B}{No se puede renunciar a ellos ni ser privado de ellos.}
  \alternativa{C}{Se adquieren al cumplir 18 años.}
  \alternativa{D}{Pueden venderse o transferirse.}
\end{frame}

\begin{frame}[allowframebreaks]{Solución}
  \framesubtitle{«Inalienable» significa que no se puede enajenar ni perder.}
  {\small
  \paso{1}{Nadie puede renunciar a sus derechos humanos ni ser despojado de ellos legítimamente.}
  }
  \alternativaok{B}{No se puede renunciar a ellos ni ser privado de ellos.}
  \begin{cajatrampa}\footnotesize
    \begin{itemize}
    \item Confundir la elección de representantes con democracia directa.
    \item Creer que los Estados deciden qué derechos humanos reconocer.
    \end{itemize}
  \end{cajatrampa}
\end{frame}

\section{Sistema económico}

\begin{frame}[plain]
  \vfill
  {\color{zinc500}\small Eje}\par\vspace{0.3em}{\Large\bfseries Sistema económico}\par\vspace{0.6em}
  {\color{zinc700} Problema económico, mercado, macroeconomía y finanzas personales.}\par\vspace{1em}
  {\small\color{zinc500} Problema económico y mercado, Macroeconomía y finanzas personales}
  \vfill
\end{frame}

\begin{frame}[allowframebreaks]{Problema económico y mercado}
  \framesubtitle{Teoría esencial}
  \begin{itemize}\small
    \item Escasez: razón de ser de la economía.
    \item Costo de oportunidad: lo que dejas de ganar al elegir.
    \item Agentes económicos: familias, empresas y Estado.
  \end{itemize}
  \begin{cajaclave}[Desplazamientos]\footnotesize
Demanda $\uparrow$ $\Rightarrow$ precio $\uparrow$ y cantidad $\uparrow$. Oferta $\uparrow$ $\Rightarrow$ precio $\downarrow$ y cantidad $\uparrow$.

Factores de la demanda: ingreso, gustos, precio de bienes relacionados. De la oferta: costos, tecnología, clima.
  \end{cajaclave}
\end{frame}

\begin{frame}{Desafío}
  \framesubtitle{Problema económico y mercado}
  {\small Una helada destruye gran parte de la cosecha de paltas. Según el modelo de oferta y demanda, ¿qué ocurrirá?}
  \vspace{0.4em}
  \alternativa{A}{El precio de la palta baja.}
  \alternativa{B}{El precio de la palta sube.}
  \alternativa{C}{La demanda de paltas aumenta.}
  \alternativa{D}{No hay cambios en el precio.}
\end{frame}

\begin{frame}[allowframebreaks]{Solución}
  \framesubtitle{La helada reduce la oferta.}
  {\small
  \paso{1}{Menos paltas disponibles: la oferta disminuye.}
  \paso{2}{Con la misma demanda, el precio de equilibrio sube.}
  }
  \alternativaok{B}{El precio de la palta sube.}
  \begin{cajatrampa}\footnotesize
    \begin{itemize}
    \item Confundir un cambio de precio (movimiento a lo largo de la curva) con un desplazamiento de la curva.
    \item Creer que lo gratuito no tiene costo de oportunidad.
    \end{itemize}
  \end{cajatrampa}
\end{frame}

\begin{frame}[allowframebreaks]{Macroeconomía y finanzas personales}
  \framesubtitle{Teoría esencial}
  \begin{itemize}\small
    \item Subir la tasa de interés frena la inflación.
    \item Inflación alta perjudica a quienes tienen ingresos fijos.
    \item La CAE permite comparar créditos con distintos plazos y cargos.
  \end{itemize}
  \begin{cajaclave}[Instrumentos]\footnotesize
Política monetaria (Banco Central): tasa de interés, meta de inflación de 3 \%. Política fiscal (Gobierno): presupuesto, gasto e impuestos.

Balanza comercial = exportaciones $-$ importaciones.
  \end{cajaclave}
\end{frame}

\begin{frame}{Desafío}
  \framesubtitle{Macroeconomía y finanzas personales}
  {\small Para controlar una inflación alta, el Banco Central de Chile suele:}
  \vspace{0.4em}
  \alternativa{A}{Bajar la tasa de política monetaria.}
  \alternativa{B}{Subir la tasa de política monetaria.}
  \alternativa{C}{Aumentar el gasto público.}
  \alternativa{D}{Fijar los precios por ley.}
\end{frame}

\begin{frame}[allowframebreaks]{Solución}
  \framesubtitle{Una tasa más alta encarece el crédito y modera el gasto.}
  {\small
  \paso{1}{Menos crédito y más ahorro reducen la demanda agregada.}
  \paso{2}{Con menor demanda, los precios suben más lento.}
  }
  \alternativaok{B}{Subir la tasa de política monetaria.}
  \begin{cajatrampa}\footnotesize
    \begin{itemize}
    \item Atribuir al Banco Central decisiones de gasto público.
    \item Confundir alza de un precio con inflación.
    \end{itemize}
  \end{cajatrampa}
\end{frame}

\section{Geografía y espacio}

\begin{frame}[plain]
  \vfill
  {\color{zinc500}\small Eje}\par\vspace{0.3em}{\Large\bfseries Geografía y espacio}\par\vspace{0.6em}
  {\color{zinc700} Población, territorio, ambiente y riesgos.}\par\vspace{1em}
  {\small\color{zinc500} Población, territorio y ambiente}
  \vfill
\end{frame}

\begin{frame}[allowframebreaks]{Población, territorio y ambiente}
  \framesubtitle{Teoría esencial}
  \begin{itemize}\small
    \item Macrocefalia: concentración de población en la capital.
    \item Amenaza natural $\neq$ desastre.
    \item Migración: factores de expulsión y de atracción.
  \end{itemize}
  \begin{cajaclave}[Riesgo]\footnotesize
\[\text{riesgo}=\text{amenaza}\times\text{vulnerabilidad}\]

Reducir la vulnerabilidad (planificación urbana, normas de construcción, educación) reduce el riesgo, aunque la amenaza no cambie.
  \end{cajaclave}
\end{frame}

\begin{frame}{Desafío}
  \framesubtitle{Población, territorio y ambiente}
  {\small ¿Qué acción reduce el riesgo de desastre ante un tsunami en una ciudad costera?}
  \vspace{0.4em}
  \alternativa{A}{Impedir que ocurran sismos submarinos.}
  \alternativa{B}{Educar a la población y señalizar vías de evacuación.}
  \alternativa{C}{Construir más viviendas en la línea de costa.}
  \alternativa{D}{Eliminar los sistemas de alerta.}
\end{frame}

\begin{frame}[allowframebreaks]{Solución}
  \framesubtitle{No se puede controlar la amenaza, pero sí la vulnerabilidad.}
  {\small
  \paso{1}{La educación y la señalización reducen la vulnerabilidad.}
  \paso{2}{Menor vulnerabilidad implica menor riesgo.}
  }
  \alternativaok{B}{Educar a la población y señalizar vías de evacuación.}
  \begin{cajatrampa}\footnotesize
    \begin{itemize}
    \item Concluir más de lo que muestran los datos.
    \item Confundir amenaza con riesgo.
    \end{itemize}
  \end{cajatrampa}
\end{frame}

\section{Habilidades de pensamiento histórico}

\begin{frame}[plain]
  \vfill
  {\color{zinc500}\small Eje}\par\vspace{0.3em}{\Large\bfseries Habilidades de pensamiento histórico}\par\vspace{0.6em}
  {\color{zinc700} Análisis de fuentes, multicausalidad, continuidad y cambio.}\par\vspace{1em}
  {\small\color{zinc500} Análisis de fuentes y pensamiento histórico}
  \vfill
\end{frame}

\begin{frame}[allowframebreaks]{Análisis de fuentes y pensamiento histórico}
  \framesubtitle{Teoría esencial}
  \begin{itemize}\small
    \item La respuesta debe sostenerse en la fuente, no en conocimientos externos.
    \item Una fuente refleja la perspectiva de su autor.
    \item Multicausalidad: evita explicaciones de una sola causa.
  \end{itemize}
  \begin{cajaclave}[Preguntas para una fuente]\footnotesize
¿Quién la produjo? ¿Cuándo y dónde? ¿Con qué propósito y para quién? ¿Qué afirma explícitamente? ¿Qué se puede inferir? ¿Qué no permite concluir?
  \end{cajaclave}
\end{frame}

\begin{frame}{Desafío}
  \framesubtitle{Análisis de fuentes y pensamiento histórico}
  {\small Un historiador analiza una carta escrita por un soldado durante la Guerra del Pacífico. ¿Qué tipo de fuente es?}
  \vspace{0.4em}
  \alternativa{A}{Secundaria, porque la analiza un historiador.}
  \alternativa{B}{Primaria, porque fue producida en la época estudiada.}
  \alternativa{C}{Secundaria, porque es un texto escrito.}
  \alternativa{D}{Primaria, porque la escribió un militar de alto rango.}
\end{frame}

\begin{frame}[allowframebreaks]{Solución}
  \framesubtitle{El criterio es la época de producción, no quién la analiza.}
  {\small
  \paso{1}{La carta se escribió durante el hecho estudiado: es primaria.}
  \paso{2}{El análisis del historiador sería una fuente secundaria.}
  }
  \alternativaok{B}{Primaria, porque fue producida en la época estudiada.}
  \begin{cajatrampa}\footnotesize
    \begin{itemize}
    \item Elegir una afirmación verdadera históricamente, pero que la fuente no respalda.
    \item Tomar la opinión del autor como un hecho comprobado.
    \end{itemize}
  \end{cajatrampa}
\end{frame}
%%</CONTENIDO>%%

\end{document}
